% Options for packages loaded elsewhere
\PassOptionsToPackage{unicode}{hyperref}
\PassOptionsToPackage{hyphens}{url}
%
\documentclass[
  11pt,
]{article}
\usepackage{lmodern}
\usepackage{amssymb,amsmath}
\usepackage{ifxetex,ifluatex}
\ifnum 0\ifxetex 1\fi\ifluatex 1\fi=0 % if pdftex
  \usepackage[T1]{fontenc}
  \usepackage[utf8]{inputenc}
  \usepackage{textcomp} % provide euro and other symbols
\else % if luatex or xetex
  \usepackage{unicode-math}
  \defaultfontfeatures{Scale=MatchLowercase}
  \defaultfontfeatures[\rmfamily]{Ligatures=TeX,Scale=1}
\fi
% Use upquote if available, for straight quotes in verbatim environments
\IfFileExists{upquote.sty}{\usepackage{upquote}}{}
\IfFileExists{microtype.sty}{% use microtype if available
  \usepackage[]{microtype}
  \UseMicrotypeSet[protrusion]{basicmath} % disable protrusion for tt fonts
}{}
\makeatletter
\@ifundefined{KOMAClassName}{% if non-KOMA class
  \IfFileExists{parskip.sty}{%
    \usepackage{parskip}
  }{% else
    \setlength{\parindent}{0pt}
    \setlength{\parskip}{6pt plus 2pt minus 1pt}}
}{% if KOMA class
  \KOMAoptions{parskip=half}}
\makeatother
\usepackage{xcolor}
\IfFileExists{xurl.sty}{\usepackage{xurl}}{} % add URL line breaks if available
\IfFileExists{bookmark.sty}{\usepackage{bookmark}}{\usepackage{hyperref}}
\hypersetup{
  hidelinks,
  pdfcreator={LaTeX via pandoc}}
\urlstyle{same} % disable monospaced font for URLs
\usepackage[margin=1in]{geometry}
\setlength{\emergencystretch}{3em} % prevent overfull lines
\providecommand{\tightlist}{%
  \setlength{\itemsep}{0pt}\setlength{\parskip}{0pt}}
\setcounter{secnumdepth}{-\maxdimen} % remove section numbering
\usepackage{mathrsfs}
\usepackage{esint}
\setmathfont{Asana Math}
\setlength{\emergencystretch}{3em}

\author{}
\date{23 September 2026}

\begin{document}

\hypertarget{a-compactified-functional-pencil-and-a-two-sided-einstein-energy-window}{%
\section{A compactified functional pencil and a two-sided Einstein
energy
window}\label{a-compactified-functional-pencil-and-a-two-sided-einstein-energy-window}}

23 September 2026. Baseline: \texttt{4401d28}.

This note extends the finite-parameter non-saturation theorem to its
vanishing-Weyl endpoint, with an explicit degenerating weight. It then
obtains the upper, as well as the lower, quantitative energy envelope.
The analytic input is the existing bounded-form compatibility theorem;
the classical feasible energy interval is not a new curvature
inequality.

\hypertarget{uniformity-on-the-unbounded-parameter-interval}{%
\subsection{1. Uniformity on the unbounded parameter
interval}\label{uniformity-on-the-unbounded-parameter-interval}}

Let \((M,g)\) be closed, connected, oriented and four-dimensional. Fix
an actual Spin\(^c\) structure whose ordinary Seiberg--Witten equations
are solvable for every metric conformal to \(g\). Write \[
 p=(a_g^+)^2>0,\quad V=\operatorname{Vol}(g),\quad
 r=\fint|\operatorname{Ric}+g|^2,\quad b_a=(-a^2)_+/V,
\] \[
 D_\beta=\frac{\beta+1}{\beta}\int(s^2+6\beta|W^+|^2)
                 -72\pi^2p,\qquad \beta\ge2.
\] The conventions are tensor Ricci norm and Hilbert--Schmidt Weyl
endomorphism norm. The class is fixed while the conformal weight varies.

\textbf{Theorem 1.} There is a universal \(c>0\) such that, for every
finite \(\beta\ge2\), \[
 \boxed{(1+b_a)(r+D_\beta/V)
       \ge c\frac{(\beta-2)^2}{\beta^4}.}                 \tag{1}
\] If \(a^2\ge T=2\chi+3\tau\), the factor \(1+b_a\) can be removed,
after reducing the universal constant. Unlike the earlier coefficient
\(c_L\), the constant in (1) does not depend on a parameter cutoff. Its
explicit weight vanishes at both the Kähler endpoint and the
vanishing-Weyl endpoint.

\textbf{Proof.} Put \(w=\sqrt6|W^+|\), \(q=-s+w\), \(e_s=s+4\),
\(z=s+\beta w\), and \(\kappa=r+D_\beta/V\). The exact nonnegative
decomposition in \href{line1_functional_pencil.md}{the finite-parameter
proof}, Section 2, implies \[
 \fint e_s^2\le4r,\quad \fint z^2\le\beta D_\beta/V,
 \quad\Delta_+/V\le D_\beta/V.
\] Consequently \[
 \fint(w-4/\beta)^2
 \le2D_\beta/(\beta V)+8r/\beta^2\le3\kappa,
\] \[
 \fint(q-q_\beta)^2\le14\kappa,
 \qquad q_\beta=4+4/\beta\in[4,6].                       \tag{2}
\] In particular the normalized Ricci, scalar and \(W^+\) squared
energies are bounded by \(C(1+\kappa)\), uniformly for \(\beta\ge2\).

Here is the required tracking of the form-production constants;
uniformity is not inferred merely by sending \(L\) to infinity. The
truncated weight \(Q=\max(q,3)\), \(f_0=\max(\sqrt Q,(-s)_+/4)\) is
independent of \(\beta\). The inherited moment argument gives the same
bounds \(|\psi|\le4\), \(D_0\le4S/27+o(1)\) and the same amplitude
control, with \(S\le117V\kappa\). Its errors involve only \(e_s\),
\(\Delta_+\) and \(Q\). For the normalized form and Weyl model the
coefficients are \[
 c_\beta^2=\frac{9\beta}{8(\beta+1)}\in[3/4,9/8],\quad
 \alpha_\beta=\frac9{8(\beta+1)}\le3/8,\quad
 t=1/\beta\in(0,1/2].
\] Set \(\omega=c_\beta\psi\) and
\(A_{\omega,t}=t(|\omega|^2I/3-\omega\otimes\omega)\). The polynomial
fit identity is \[
 |W^+-A_{\omega,t}|^2
 =\tfrac16(w-2\alpha_\beta|\psi|^2)^2
                                  +2\alpha_\beta E_\lambda.
\] The amplitude estimate, (2), and
\(2\alpha_\beta|\psi|^2=u^2/(\beta+1)\) therefore yield, with one smooth
actual solution and an arbitrarily small allowance, \[
 \|\omega\|_\infty\le3\sqrt2,\quad
 d_t:=\fint\bigl(|\operatorname{Ric}+g|^2
   +|W^+-A_{\omega,t}|^2+(|\omega|^2-2)^2\bigr)
                       \le C\kappa+o(1).                 \tag{3}
\] No inverse power of \(t\) is used here. The determinant energy
identity and the bounded-spinor Hessian estimate then give \[
 B_\omega:=\fint|\Delta_H\omega|^2
                   \le C(1+\kappa+b_a)+o(1),             \tag{4}
\] again uniformly: the needed curvature energies are bounded by (2),
and the normalization of the quadratic spinor map is bounded above. This
is exactly the estimate proved in Sections 3--5 of the finite-parameter
paper, with every parameter-dependent coefficient now bounded
explicitly.

The bounded-form compatibility theorem, with the fixed parameter range
\(0\le t\le1/2\), gives \[
 t^2(1-2t)^2\le C\kappa(1+\kappa+b_a).                   \tag{5}
\] For \(\kappa\le1\) this proves (1); for \(\kappa>1\) use
\(t^2(1-2t)^2\le1/64\). Let the smoothing allowance tend to zero. If
\(a^2\ge T\), Chern--Weil gives \(b_a\le r/(8\pi^2)\), and the same
split removes the square-cost factor. This proves the theorem. No
derivative estimate for the harmonic projection is asserted uniformly as
\(t\to0\). \(\square\)

\hypertarget{a-pencil-with-an-actual-endpoint-at-zero}{%
\subsection{2. A pencil with an actual endpoint at
zero}\label{a-pencil-with-an-actual-endpoint-at-zero}}

Define on \(0\le t\le1/2\) \[
 P_t=t(1+t)\int s^2+6(1+t)\int|W^+|^2-72\pi^2tp.
                                                               \tag{6}
\] For \(t>0\), this is \(tD_{1/t}\); at zero it is exactly
\(P_0=6\int|W^+|^2\). Hence \(P_t\ge0\) on the closed interval.
Multiplying (1) by \(t\), and using \(r+tD/V\ge t(r+D/V)\), proves \[
 \boxed{(1+b_a)(r+P_t/V)\ge c t^3(1-2t)^2.}             \tag{7}
\] The same square-cost removal applies when \(a^2\ge T\). At \(t=0\)
the inequality is trivial and supplies no nonzero lower amplitude or
actual anti-self-dual Einstein existence. At \(t=1/2\), the endpoint
includes negative Kähler--Einstein metrics. The cubic decay in (7) is a
proved sufficient rate from this transfer, not an optimality claim.

For a normalized negative Kähler--Einstein metric and its canonical
class, \[
 P_t/V=4(1-2t)^2.
\] This provides an actual normalization check. It does not establish
sharpness at the other endpoint.

\hypertarget{the-missing-upper-energy-envelope}{%
\subsection{3. The missing upper energy
envelope}\label{the-missing-upper-energy-envelope}}

Assume now that \(A>0\), \[
 T=2\chi+3\tau=A(1-h),\qquad 1/4<h<1/3,                 \tag{8}
\] and for every metric one can choose an actual conformally solvable
class with \(p\ge A\) and \(a^2\ge T\). Put \[
 e=\frac{\int|\operatorname{Ric}|^2}{8\pi^2A},\quad
 F(h)=\frac{(4-\sqrt{1-3h})^2}{9},\quad
 G(h)=\frac{(4+\sqrt{1-3h})^2}{9},
\] \[
 t_*(h)=\frac{1-2\sqrt{1-3h}}{8+2\sqrt{1-3h}}\in(0,1/8),
 \qquad \Theta(h)=t_*^3(1-2t_*)^2.                       \tag{9}
\] Use the all-sign relative defect \[
 \mathcal E(g)=\inf_{\lambda\ne0}
 \frac{\int|\operatorname{Ric}-\lambda g|^2}{\lambda^2V}.
\]

\textbf{Theorem 2.} There is a universal \(c_*>0\) such that every
smooth metric under (8) satisfies \[
 \boxed{\mathcal E(g)+(G(h)-e)_+\ge c_*\Theta(h).}        \tag{10}
\] Combining with the inherited lower-envelope theorem, every Einstein
candidate must satisfy the strictly smaller energy window \[
 \boxed{F(h)+\varepsilon h^2\le e
                       \le G(h)-c_*\Theta(h).}          \tag{11}
\] Ricci-flat and positive Einstein metrics are independently excluded
by the positive actual monopole scalar estimate. Equation (11) concerns
all remaining Einstein metrics, not just minimizers or a gluing
neighborhood.

\textbf{Proof.} Write \(R=\int|\operatorname{Ric}|^2\),
\(P=\int|\operatorname{Ric}^{\circ}|^2\), \(C_A=8\pi^2A\). The exact
Chern--Weil reduction of (6) is \[
 P_t=k(t)R+\ell(t)P+12\pi^2(1+t)T-72\pi^2tp,
\] \[
 k(t)=\frac{(1+t)(8t-1)}2,\qquad
 \ell(t)=2(1+t)(1-2t).                                  \tag{12}
\] For \(t=t_*(h)\), direct substitution gives \[
 k(t)G(h)+\tfrac32(1+t)(1-h)-9t=0,
\] and therefore \[
 P_t=C_A k(t)(e-G(h))+\ell(t)P-72\pi^2t(p-A)
       \le C_A|k(t)|(G(h)-e)_++\ell(t)P.                \tag{13}
\] This is the opposite sign of the coefficient in the lower-envelope
argument. It is why the finite-parameter range beyond \(\beta=8\)
controls an upper energy wall.

First fix the target to \(\operatorname{Ric}=-g\) and
\(r=\fint|\operatorname{Ric}+g|^2\le1\). The scalar monopole bound and
the trace estimate imply \[
 C_A/V\le(2+\sqrt r)^2\le9,\qquad P/V\le r.
\] On \(0\le t\le1/8\), \(|k(t)|\le1/2\) and \(0\le\ell(t)\le2\). Using
(7), with its square-cost factor removed, in (13) gives \[
 c\Theta(h)\le3r+\tfrac92(G(h)-e)_+
              \le\tfrac92\{r+(G(h)-e)_+\}.             \tag{14}
\] For \(r>1\) the desired estimate follows by decreasing the constant.
Constant rescaling and taking the infimum over negative target constants
prove (10) for the negative-target infimum \(\mathcal E_-\).

The passage to all signs needs a different argument from the earlier
lower-envelope transfer: a very large \(e\) now makes the energy excess
zero. Here the upper Chern--Weil budget closes that possibility. Let
\(I=\int s\). If \(I\le0\), direct minimization gives
\(\mathcal E_+=4\). If \(I>0\), it gives
\(\mathcal E_+=\mathcal E=4-I^2/(RV)\). The signed scalar estimate
\(\int s_-^2\ge4C_A\) implies \[
 \mathcal E_+\ge4/e.                                   \tag{15}
\] If \(\mathcal E_+<1/2\), the chirality upper bound gives \[
 e(1-\mathcal E_+)\le
 U=3\min(2\chi+3\tau,2\chi-3\tau)/A\le3,
\] so \(e\le6\). Equation (15) would imply \(\mathcal E_+\ge2/3\), a
contradiction. Thus \(\mathcal E_+\ge1/2\). Since \(\Theta\) is bounded,
this proves (10) also for the positive target infimum, after reducing
\(c_*\); take the minimum of the two signwise bounds. \(\square\)

\hypertarget{geometry-of-the-two-endpoints-and-an-actual-application}{%
\subsection{4. Geometry of the two endpoints and an actual
application}\label{geometry-of-the-two-endpoints-and-an-actual-application}}

The classical Einstein constraints are, with nonnegative \(x,y\), \[
 x^2=4e,\quad x^2+8y^2=12(1-h),\quad x\ge2,\quad x+y\ge3.
\] When \(1/4<h<1/3\), the ellipse has \(x<3\). Solving these
inequalities gives the full classical interval \(F(h)\le e\le G(h)\).
Thus the interval itself is classical; the positive quantitative margins
in (11) are the new application of the differential non-saturation
mechanism.

At \(h\downarrow1/4\), the upper optimal parameter tends to \(t=0\), and
\(\Theta(h)\sim(8/27)(h-1/4)^3\). At \(h\uparrow1/3\), both optimal
parameters meet \(t=1/8\); the interval collapses, consistently with the
old critical gap. At the lower-envelope endpoint \(h\downarrow0\), the
optimal parameter instead tends to the allowed Kähler value \(t=1/2\).
These are different geometric degeneracies; they cannot be treated as
one compact parameter range with a fixed forbidden distance.

For an actual example, let \[
 M=(\Sigma_2\times\Sigma_3)\#5\overline{\mathbb{CP}}^{\,2}.
\] Its polarized actual classes give \(A=16\), \(h=5/16\), \(a^2=T=11\).
Here \[
 F=25/16,\qquad G=289/144,\qquad t_*=1/17,
 \qquad \Theta=225/17^5.
\] Any Einstein metric, if one exists, must consequently satisfy \[
 \frac{25}{16}+\varepsilon\frac{25}{256}
 \le\frac{R}{128\pi^2}
 \le\frac{289}{144}-c_*\frac{225}{17^5}.                 \tag{16}
\] This is not a nonexistence claim for this example. Its bare
two-chirality upper bound is only \(e\le33/16\), so the upper envelope
contains information not supplied by that bound alone. The actual
monopole selection is the same established polarized input as in the
main paper; no new gauge existence theorem is asserted here.

\hypertarget{relation-to-the-research-goal}{%
\subsection{5. Relation to the research
goal}\label{relation-to-the-research-goal}}

The first-line theorem now supplies a two-sided restriction on the
entire energy spectrum of possible Einstein outputs in the displayed
regime. It does not merely study an energy-minimizing sequence. For a
pointwise input, \(\mathcal E\le4\delta_\sigma^2\) converts (10) into
the corresponding upper-energy constraint on that input; a small
residual alone does not force the input energy into a forbidden region.

This note neither classifies the allowed spectra of all Einstein metrics
nor constructs an Einstein metric. The separate Bauer--Furuta
application in this increment is the new topological-category result.
The present contribution is the closed-parameter formulation and its
upper-wall application, using the existing analytic mechanism.

Sources: \href{https://arxiv.org/abs/math/0003068}{LeBrun's scalar--Weyl
inequalities and Chern--Weil reduction}; the local
\href{line1_functional_pencil.md}{functional-pencil proof}, Sections
2--6; \href{round21_analytic.md}{bounded-form compatibility}, Section 6;
\href{line1_energy_envelope_gap.md}{lower energy envelope}; and
\href{line1_two_chirality_topology.md}{two-chirality upper bound}.

\texttt{check\_line1\_compactified\_pencil.py} verifies the exact
coefficients, both envelope roots, endpoint rates, actual example, and
several deliberate sign errors. These checks do not substitute for the
written uniform analytic estimates or actual gauge-solvability inputs.

\end{document}
